\documentclass[11pt,a4paper]{article}
\usepackage[margin=18mm,top=17mm,bottom=19mm]{geometry}
\usepackage{amsmath,amssymb,graphicx,array,fancyhdr,lastpage,textcomp}
\usepackage{fontspec}
\setmainfont{texgyreheros-regular.otf}[BoldFont=texgyreheros-bold.otf,ItalicFont=texgyreheros-italic.otf,BoldItalicFont=texgyreheros-bolditalic.otf]
\setlength{\parindent}{0pt}
\setlength{\parskip}{3pt}
\pagestyle{fancy}\fancyhf{}
\renewcommand{\headrulewidth}{0pt}\renewcommand{\footrulewidth}{0.3pt}
\fancyfoot[L]{\footnotesize by paper.shrit.in\quad |\quad normal}
\fancyfoot[R]{\footnotesize Page \thepage\ of \pageref{LastPage}}
\newcommand{\question}[3]{\par\vspace{8pt}\noindent\begin{minipage}[t]{0.075\linewidth}\textbf{#1)}\end{minipage}\begin{minipage}[t]{0.855\linewidth}#3\end{minipage}\hfill\begin{minipage}[t]{0.055\linewidth}\raggedleft(#2)\end{minipage}\par}
\begin{document}
{\LARGE\bfseries Question Paper}\hfill {\small Registration No.: \rule{33mm}{0.3pt}}\par\vspace{8pt}
\begin{center}{\large\bfseries MANIPAL FORMAT | PRACTICE PAPER}\\[6pt]{\bfseries THIRD SEMESTER B.TECH. | MIDSEM PRACTICE}\\[4pt]{\bfseries DATA COMMUNICATION AND COMPUTER NETWORKS [CSS 2102]}\\[4pt]{\small COMPUTER SCIENCE AND ENGINEERING}\\[3pt]{\small SET B: NEWLY AUTHORED QUESTIONS}\end{center}
\textbf{Marks: 30}\hfill\textbf{Suggested duration: 90 minutes}\par\vspace{3pt}\hrule\vspace{5pt}
{\small\textbf{Answer all questions.} Questions 1--5 are MCQs worth 1 mark each. Select one correct option per MCQ. The remaining questions are theory questions worth 25 marks in total; show working and draw labelled diagrams where appropriate. Modules 1-2. Independent practice paper; not an official university examination.}\par
\medskip\textbf{Section A: Multiple-choice questions (5 marks)}\par
\question{1}{1}{A noiseless channel has bandwidth 3 kHz and four signal levels. What is its Nyquist maximum bit rate?\par (A) 3 kbps\par (B) 6 kbps\par (C) 12 kbps\par (D) 24 kbps}
\question{2}{1}{How many bits can one of four distinct signal levels represent?\par (A) 1\par (B) 2\par (C) 4\par (D) 8}
\question{3}{1}{A channel has bandwidth 3 kHz and linear signal-to-noise ratio 31. What is its Shannon capacity?\par (A) 15 kbps\par (B) 12 kbps\par (C) 93 kbps\par (D) 30 kbps}
\question{4}{1}{For the same channel and signalling scheme, Nyquist gives 12 kbps and Shannon gives 15 kbps. Which is the tighter upper bound?\par (A) 15 kbps\par (B) 27 kbps\par (C) 3 kbps\par (D) 12 kbps}
\question{5}{1}{Under the noiseless Nyquist model, what happens to the maximum bit rate when bandwidth doubles and the number of levels stays fixed?\par (A) It halves\par (B) It stays fixed\par (C) It doubles\par (D) It quadruples}
\newpage\textbf{Section B: Theory questions (25 marks)}\par
\question{6}{3}{Compare ASK, binary FSK and BPSK in terms of the carrier parameter changed and sensitivity to amplitude noise. State one reason to prefer BPSK over ASK in a channel with amplitude disturbances.}
\question{7}{2}{An application generates 900 payload bytes. Each of five layers adds a 20-byte header. Calculate the total transmitted size and the percentage occupied by headers, ignoring trailers.}
\question{8}{5}{Encode data \texttt{1101011011} using CRC generator \texttt{10011}. Show modulo-2 division, the four-bit remainder and the complete transmitted codeword. State the receiver\textquotesingle s check for an error-free received word.}
\question{9}{3}{Draw Manchester and bipolar AMI waveforms for \texttt{10110010}. For Manchester, use low-to-high at mid-bit for 1 and high-to-low for 0. For AMI, assume the last nonzero pulse before this sequence was negative. Label all bit boundaries.}
\question{10}{2}{Explain the different actions of a hub and a learning switch when a frame arrives. For the switch, distinguish known and unknown destination MAC addresses.}
\question{11}{5}{A stop-and-wait link sends 1000-byte data frames at 1 Mbps. One-way propagation delay is 10 ms. ACK transmission and processing times are negligible; no frames are lost. Calculate frame transmission time, time per successful cycle, utilization and useful throughput. Ignore headers.}
\question{12}{3}{A selective-repeat sender sends frames 0, 1, 2 and 3 with window size 4. Frame 1 is lost; frames 0, 2 and 3 arrive and their individual ACKs return before timeout. No new frames are available. State which frames are buffered, which frame is retransmitted and the final in-order delivery sequence. Compare with go-back-N under the same loss, assuming it discards out-of-order frames.}
\question{13}{2}{A shared half-duplex Ethernet link operates at 100 Mbps. The maximum one-way propagation delay is $5\,\mu$s. Ignoring other delays, calculate the minimum frame length required for collision detection in bits and bytes.}
\vfill\begin{center}\small End of question paper\end{center}\end{document}